\documentclass{resume}

\begin{document}

\fontfamily{ppl}\selectfont

\noindent
\begin{tabularx}{\linewidth}{@{}m{0.8\textwidth} m{0.2\textwidth}@{}}
{
    \Large{Luciano Ortiz} \newline
    \small{
        \clink{
            \href{mailto:lucianortiz97@gmail.com}{lucianortiz97@gmail.com} \textbf{|} 
            {\fontdimen2\font=0.75ex +11 6657 8365} 
            \textbf{|} \href{https://www.linkedin.com/in/luciano-ortiz-521bb3117/}{linkedin} \textbf{|}
            \href{https://github.com/luctiz}{github} 
            \textbf{|}
             
        } \newline
        Buenos Aires, Argentina
    }
}
\end{tabularx}
\begin{center}
\begin{tabularx}{\linewidth}{@{}*{2}{X}@{}}
% left side %
{   \csection{DESCRIPCIÓN}{\bigskip \small


            Desarrollador de software con más de 3 años de experiencia en desarrollo web full stack, especializado en .NET. Tengo experiencia en metodologías ágiles y puedo mantener comunicaciones fluidas en inglés, tanto escrito como oral. Mis habilidades abarcan desde la creación de interfaces frontend responsive adaptadas a los requerimientos de UX/UI, hasta la implementación de backends con arquitectura en capas, API REST, ORM y tests unitarios.
            \newline
            Además soy aficionado de la tecnología y la programación, con muchas ganas de seguir aprendiendo y desempeñarme lo mejor posible en este ámbito.
            \newline
            }

                    \csection{Experiencia laboral}{\small
        \begin{itemize}
            % item 2 %
            \item \frcontent{SMS Sudamérica - Desarrollador Fullstack Ssr.}{Trabajo como desarrollador fullstack utilizando metodologías ágiles de trabajo. Utilizo principalmente Angular para el frontend y .NET6/API Controllers para el backend. La mayoría de las conversaciones las mantengo en inglés, tanto escrito como oral}
            {.NET, Entity Framework, Angular, React, Apache Kafka, Node.JS, Docker, SQL Server, Flyway}
            {Julio 2023 - Agosto 2024}
            \item \frcontent{AppDirect - Fullstack Developer Engineer}{Trabajo como desarrollador fullstack utilizando metodologías ágiles de trabajo. Utilizo principalmente Angular para el frontend y .NET6/API Controllers para el backend. La mayoría de las conversaciones las mantengo en inglés, tanto escrito como oral}
            {.NET, Entity Framework, Angular, React, Apache Kafka, Node.JS, Docker, SQL Server, Flyway}
            {Julio 2023 - Agosto 2024}
            \item \frcontent{Axton IT - Desarrollador Fullstack}{Desarrollé soluciones web nuevas y modifiqué existentes de manera de cumplir los requerimientos de los clientes de la empresa. 
            Utilicé .Net y .Net Core, con C\#, Visualbasic, Entity Framework, y Blazor con MVC Controllers para backend. Para el desarrollo frontend trabajé con HTML/CSS/JS puro o React/Blazor. Utilicé diversos tipos de arquitecturas como cliente/servidor, API, MVC, entre otras. \newline Además utilicé el framework Flask/Python para armar y desplegar un algoritmo de inteligencia artificial para un chatbot de la empresa.}{}{Abril 2021 - Junio 2023}

        \end{itemize}
} 
% end left side %
& 
% right side %
{

    \csection{EDUCACIÓN}{\small
        \begin{itemize}
            % item 1 %
            \item \frcontent{Lic. En Analista de Sistemas}{Universidad de Buenos Aires}{}{2015 - Actualidad}
        \end{itemize}
    }
    
    \csection{HABILIDADES}{\small
        \begin{itemize}
            \item \textbf{Tecnologías} \newline
            {\footnotesize  .NET, C\#, C++, HTML/CSS, Javascript, Blazor/Razor, React, Python, scikit-learn, Entity Framework, SQL, Microsoft SQL Server, WebGL/OpenGL, Java, Git, Heroku}
             \item \textbf{Idiomas} \newline
           {\footnotesize  Español, Inglés B2 }
        \end{itemize}
    }
        
    }
    
    \csection{Proyectos universitarios}{\small
        \begin{itemize}
            \item \frcontent{Wolfenstein 3D online 
            \clink{\href{https://github.com/German-Rotili/tpGrupal}{[código]}}}
            {Un remake del shooter Wolfenstein 3D con soporte para multijugador online mediante arquitectura cliente-servidor}{}{C++}
            \item \frcontent{Proyecto ITIL  
            \clink{\href{https://itil-front.herokuapp.com/}{[demo]}} 
            \clink{\href{https://github.com/AndresJalife/itil-front}{[código front-end]}}}
            {Herramienta de gestión de incidentes utilizando las buenas prácticas ITIL. Junto a mi equipo, me encargué principalmente de programar el front utilizando React}{}{ React}
        \end{itemize}
    }
    \csection{Otros aspectos destacados}{\small
        \begin{itemize}
            \item {\footnotesize Conocimientos de machine learning / análisis de datos / big data, aprendidos tanto en la universidad como en el trabajo}
            \item {\footnotesize Conocimientos de las bases de Computación gráfica. Llevados a la práctica utilizando WebGL.
            \clink{\href{https://github.com/luctiz/sg-tp}{[codigo]}}}
            \clink{\href{https://luctiz.github.io/sg-tp/}{[demo]}}
        \end{itemize}
    }
  
}
\end{tabularx}
\end{center}
\end{document}